\section{Background: the geometric product and angular momentum}\label{sec:background}

\textbf{Multivectors.} The algebra $\mathrm{Cl}(3,0)$ is generated by orthonormal basis vectors $e_1,e_2,e_3$ with $e_i^2=1$ and $e_ie_j=-e_je_i$ for $i\neq j$ \cite{lounesto2001clifford}. Its $2^3=8$ basis elements are graded by the number of vectors multiplied together: one scalar (grade 0), three vectors $e_i$ (grade 1), three bivectors $e_ie_j$ representing oriented planes (grade 2), and one pseudoscalar $e_1e_2e_3$ representing an oriented volume (grade 3), with a multivector being a linear combination of all eight \cite{ruhe2023gcan}.

\textbf{The geometric product.} For two vectors $u,v$ the geometric product splits into a symmetric and an antisymmetric part \cite{brandstetter2022clifford},
\begin{equation}
uv = u\cdot v + u\wedge v ,
\label{eq:gp}
\end{equation}
a scalar plus a bivector. In three dimensions the Hodge dual maps the bivector $u\wedge v$ to the vector $u\times v$ \cite{doran2003geometric}. The same product applies uniformly to multivectors of any grade and commutes with rotations, enabling equivariant Clifford network operations \cite{ruhe2023cgenn}.

\textbf{Angular momentum.} An irreducible $\mathrm{SO}(3)$ feature of order $L$ has $2L{+}1$ components: scalars are $L{=}0$, vectors are $L{=}1$, and symmetric traceless $3\times3$ matrices are $L{=}2$ \cite{varshalovich1988quantum}. The outer product of two vectors has the classical Cartesian decomposition \cite{coope1965irreducible}
\begin{equation}
u_iv_j = \underbrace{\tfrac13\delta_{ij}\,u\cdot v}_{L=0} + \underbrace{\tfrac12(u_iv_j-u_jv_i)}_{L=1} + \underbrace{\tfrac12(u_iv_j+u_jv_i)-\tfrac13\delta_{ij}\,u\cdot v}_{L=2} .
\label{eq:decomp}
\end{equation}
Comparing eqn~(\ref{eq:gp}) and (\ref{eq:decomp}): the scalar $u\cdot v$ is the $L{=}0$ term and the bivector $u\wedge v$ is the $L{=}1$ term, but the geometric product omits the symmetric traceless $L{=}2$ matrix component used in Cartesian tensor representations \cite{simeon2023tensornet}. Bivectors carry the axial $L{=}1$ representation despite having grade 2 \cite{geiger2022e3nn}; applying the scalar, vector and dual transformation rules to the graded components gives
\begin{equation}
\mathrm{Cl}(3,0) \cong 2\cdot(L{=}0) \oplus 2\cdot(L{=}1), \qquad L{=}2 \notin \mathrm{Cl}(3,0) .
\label{eq:cl_irreps}
\end{equation}
The $L{=}2$ term of eqn~(\ref{eq:decomp}) is exactly the symmetric traceless projection that CliffordSTF adds in closed form. A fuller account of the algebra is given in ESI Section~\ref{E-supp:prelim}.
